We retain the notation of Section~\ref{sec:unitary-geometry} and
write $\mu=\lambda(x)$. The two rank-one families are the
decomposable and stable-system families of \cref{m2:rat:rankone-param}.

The following rational examples illustrate the two rank-one
families. We set
\[
 P_\pm=\frac12\begin{pmatrix}1&\pm1\\\pm1&1\end{pmatrix},
 \qquad
 x_{\mathrm{dec}}=\begin{pmatrix}P_-&P_+\\P_+&P_-\end{pmatrix}.
\]
The two projectors are $J$-selfadjoint and have orthogonal
images.  The matrix swaps the two $P_+$ lines and fixes
the two $P_-$ lines, so it is unitary.  It is a
decomposable point.  For $t\in F^\times$, $t^2\ne1$,
let $h_t=\diag(t,t^{-1})$ and
\begin{equation}\label{m2:rat:rankone-examples}
 x_t=x_{\mathrm{dec}}\diag(I,h_t).
\end{equation}
This is again unitary.  Its incidence lines are
$\im P_+$ and $h_t^{-1}\im P_-$, neither stable under
$D=P_-h_t$, so it is a stable system.  Both examples have multiplier one.

We use the real-pair coordinates of \cref{m2:rat:real-pair}.

\begin{lemma}\label{m2:rat:dense-transporter}
For every dense closed point
$x=\left(\begin{smallmatrix}A_x&B_x\\C_x&D\end{smallmatrix}\right)$
there is a rational transporter $k_x=(u_x,v_x,h_x)\in K(F)$
with $k_x\cdot x^t=x$.  It can be constructed by
linear equations over $F$.  If $D$ is scalar, one may
choose $h_x=1$.
\end{lemma}

\begin{proof}
In the real-pair coordinates of \cref{m2:rat:real-pair},
closedness implies that
$(A,B_0)$ is geometrically semisimple: its generated
algebra after extension to $E$, and hence to $\barF$,
is the algebra in \eqref{m2:rat:dense-algebra}.
Choose an invertible symmetric solution $H$ over $F$
of
\[
 AH=HA^t,\qquad B_0H=HB_0^t
\]
by \cref{m2:geo:pair-transpose}.  This is a linear system
on three entries; the lemma shows that an invertible
solution exists in every semisimple case.  Put
\begin{equation}\label{m2:rat:dense-h}
 h_x=P(HR^{-1})P^{-1}.
\end{equation}
By \eqref{m2:rat:real-transpose},
$h_xD^th_x^{-1}=D$, and by
\cref{m2:rat:real-pair} $h_x\in G_2(F)$.
Moreover $R^{-1}=P^2$, so $h_x=PHP$ is symmetric.
Define
\begin{equation}\label{m2:rat:dense-k}
 u_x=B_xh_xC_x^{-t},\qquad
 v_x=C_x^{-1}h_xB_x^t.
\end{equation}
The $B,C,D$ transporter equations follow immediately.
We check the similitudes as well.  Writing
$\lambda=\lam(h_x)$, the relation
$Dh_x=h_xD^t$ gives
\[
 h_x^\dagger Q_Dh_x
 =\lambda(\mu J-\bar D J D^t)
 =\lambda\bar C_x J C_x^t,
\]
where the last equality is the transpose of the last
row equation in \eqref{m2:geo:row-equations}.  Therefore
$u_x^\dagger Ju_x=\lambda J$.
Symmetry of $h_x$ gives
$v_x=C_x^{-1}h_x^tB_x^t=u_x^t$.
Ordinary transpose preserves $G_2$ and its multiplier,
so $v_x^\dagger Jv_x=\lambda J$ as well.
Thus $(u_x,v_x,h_x)\in K(F)$.
Acting by this element on $x^t$ produces a unitary
matrix with the same multiplier and the same $B,C,D$
blocks as $x$.  By the uniqueness of $A=B_xX_DC_x$
in \cref{m2:rat:dense-param}, its $A$-block also agrees.

If $D$ is scalar, it is fixed by transpose, and one
can take $h_x=1$ directly in \eqref{m2:rat:dense-k}.
Here $B_x$ and $C_x$ have the same multiplier
$\mu-\N(D)$, by the two Gram equations, so the two
remaining factors have multiplier one.  The same
uniqueness argument verifies the $A$-block.
\end{proof}

For rank zero write $P^{-1}DP=A+\delta B_0$ over $F$
as in \cref{m2:rat:real-pair}.  Since $D$ is semisimple
and $D^*=\mu D^{-1}$, the algebra
$E\langle D,D^*\rangle=E[D]$ is semisimple.
Thus the real pair $(A,B_0)$ is geometrically semisimple.
Applying \cref{m2:geo:pair-transpose} to this pair and
using \eqref{m2:rat:dense-h} gives $h_x\in G_2(F)$ with
$h_xD^th_x^{-1}=D$.  Choose any $v_x\in G_2(F)$
of multiplier $\lam(h_x)$, for example
$\diag(1,\lam(h_x))$, and set
$u_x=a_\mu v_xa_\mu^{-1}$.  This gives an explicit
transporter for \eqref{m2:rat:rankzero-representative}.

For rank one we use the scalar endomorphism algebra to construct
a transporter.

\begin{proposition}\label{m2:rat:rankone-transporter}
Let $x$ be an arbitrary rank-one closed $F$-point.
A rational transporter from $x^t$ to $x$ can be
constructed by solving an $E$-linear system followed
by quadratic local norm equations.
\end{proposition}

\begin{proof}
First solve, without imposing the unitary condition,
\begin{equation}\label{m2:rat:linear-transporter}
 u_0A^t=Av_0,\quad u_0C^t=Bh_0,\quad
 h_0B^t=Cv_0,\quad h_0D^t=Dh_0.
\end{equation}
These equations are linear in the entries of
$u_0,v_0,h_0$ over $E$.  An invertible solution exists
by \cref{m2:geo:transpose} applied to the split-quiver
representation over $E$.  After row reduction it may
be selected by avoiding the zero locus of the product
of the three determinants in the free variables.
Let $t_0=(u_0,v_0,h_0)$ be such a solution, let
$y=x^t$, and write $L_0y=xR_0$ for its block equation.

Set
\[
 a=(u_0^*u_0,v_0^*v_0,h_0^*h_0).
\]
The unitary equations for $x,y$, with the same
multiplier $\mu$, imply $L_0^*x=yR_0^*$.
It follows that
\[
 L_0^*L_0y=yR_0^*R_0.
\]
Thus $a$ is an invertible selfadjoint element of
$\mathscr A_y=E^m$, with $m=3$ or $2$.
Its decomposition along the selfadjoint primitive idempotents
of \cref{m2:rat:selfadjoint-idempotents} is
$a=\sum_i a_ie_i$, with $a_i\in F^\times$.

At every vertex the determinant of $t_0^*t_0$ is
$\N_{E/F}(\det t_0)$, hence a norm.  In the
decomposable case the first and middle vertices
therefore say that $a_ra_a$ and $a_ra_b$ are norms.
In the stable-system case the first vertex says
that $a_ra_a$ is a norm.  Since $Q_F$ has exponent
two, in either case all the $a_i$ have the same
norm class, by the incidence-determinant computation of
\cref{m2:rat:rankone-kernel} applied to $t_0^*t_0$.
Set $\lambda=a_r$ and choose
$s_i\in E^\times$ with
\[
 \N(s_i)=\lambda/a_i.
\]
These are solvable local norm equations.  For
$s=\sum_i s_ie_i\in\mathscr A_y^\times$ one has
$s^*as=\lambda I$ at every vertex.  Hence
$k_x=t_0s$ is a triple of Hermitian similitudes
with the same multiplier $\lambda$.  Since $s$
stabilizes $y$, the triple $k_x$ still transports
$y$ to $x$.  This constructs the required element
of $K(F)$.
\end{proof}

For the examples at the start of this appendix, $x_{\mathrm{dec}}$ is
symmetric and one may take $k_x=1$.  For
\eqref{m2:rat:rankone-examples}, direct block
multiplication gives the transporter
$(I,I,h_t^{-1})$.

